\phantomsection
\section*{Введение}
\addcontentsline{toc}{section}{Введение}

\textbf{Актуальность темы.} Рынок персональных компьютеров активно развивается, а спрос на индивидуальные сборки ПК стабильно сохраняется: пользователи заинтересованы в получении устройств, максимально соответствующих их требованиям по производительности, стоимости, внешнему виду и назначению. Это приводит к росту числа небольших компаний и частных специалистов, занимающихся продажей комплектующих, сборкой компьютеров под заказ и сервисным сопровождением. Для таких организаций критически важна автоматизация процессов учёта комплектующих, формирования сборок, оформления заказов и анализа результатов продаж: при отсутствии специализированной информационной системы возникают проблемы с контролем остатков на складе, отслеживанием движения товаров, ведением клиентских заказов и расчётом прибыли.

\textbf{Степень разработанности проблемы.} На рынке представлено множество интернет-магазинов компьютерной техники (DNS, «Регард», «НИКС», «Ситилинк») и универсальных CRM-/ERP-систем. Однако существующие решения либо ориентированы на крупный розничный сегмент и не учитывают специфику внутреннего учёта комплектующих при сборке ПК под заказ, либо избыточны и дороги для малого бизнеса. Это обуславливает целесообразность разработки специализированного веб-приложения, объединяющего функции каталога, склада, конфигуратора, заказов и аналитики в едином интерфейсе с разграничением прав доступа.

\textbf{Объект исследования} --- процессы учёта комплектующих, сборки и продажи персональных компьютеров в малых компаниях.

\textbf{Предмет исследования} --- веб-приложение, предназначенное для автоматизации учёта комплектующих, управления сборками ПК и оформления заказов с разграничением прав доступа по ролям.

\textbf{Целью данной работы} является проектирование и реализация веб-приложения, обеспечивающего автоматизацию учёта комплектующих, управление сборками ПК и взаимодействие администраторов, менеджеров и клиентов через систему личных кабинетов.

\textbf{Для достижения цели} необходимо решить следующие задачи:
\begin{itemize}
    \item провести анализ предметной области и обзор существующих программных решений;
    \item обосновать выбор технологического стека для разработки;
    \item спроектировать архитектуру приложения, модель данных и пользовательские интерфейсы;
    \item реализовать аутентификацию, ролевую модель и CRUD-операции над основными сущностями;
    \item реализовать конфигуратор сборки, аналитические отчёты, импорт и экспорт данных;
    \item выполнить развёртывание приложения и оформить пояснительную записку.
\end{itemize}

\textbf{Практическая значимость} работы заключается в создании готового к внедрению программного продукта, который может использоваться действующими магазинами комплектующих и небольшими сборщиками ПК для автоматизации ключевых бизнес-процессов и сокращения операционных издержек.
